\documentclass[10pt,a4paper]{amsart}
\usepackage{amsmath,amssymb,amsthm,graphicx}
% Retained minimal preamble for the Stacks Project chapter more-algebra.tex (More on Algebra).
% The chapter file pinned-original-source.tex carries only \input{preamble} on its first line, so no
% preamble text lives inside the chapter itself.  The Stacks Project uses one single shared preamble for
% every chapter; the snapshot stages it identically as sources/stacks-chow/preamble.tex and
% sources/stacks-etale-cohomology/preamble.tex.  The block below is copied mechanically, line for line,
% from sources/stacks-chow/preamble.tex lines 14-18 (the Xy-pic setup) and 154-217 (the native theorem
% declarations together with the mathematical macros).  Nothing was added, removed or redefined; the
% wrapper class and the amsmath/amssymb/amsthm/graphicx base packages are supplied by the assembler.
% Chapter: more-algebra.tex -- More on Algebra.
% For commutative diagrams we use Xy-pic
\usepackage[all]{xy}

% We use 2cell for 2-commutative diagrams.
\xyoption{2cell}

% Theorem environments.
%
\theoremstyle{plain}
\newtheorem{theorem}[subsection]{Theorem}
\newtheorem{proposition}[subsection]{Proposition}
\newtheorem{lemma}[subsection]{Lemma}

\theoremstyle{definition}
\newtheorem{definition}[subsection]{Definition}
\newtheorem{example}[subsection]{Example}
\newtheorem{exercise}[subsection]{Exercise}
\newtheorem{situation}[subsection]{Situation}

\theoremstyle{remark}
\newtheorem{remark}[subsection]{Remark}
\newtheorem{remarks}[subsection]{Remarks}

\numberwithin{equation}{subsection}

% Macros
%
\def\lim{\mathop{\mathrm{lim}}\nolimits}
\def\colim{\mathop{\mathrm{colim}}\nolimits}
\def\Spec{\mathop{\mathrm{Spec}}}
\def\Hom{\mathop{\mathrm{Hom}}\nolimits}
\def\Ext{\mathop{\mathrm{Ext}}\nolimits}
\def\SheafHom{\mathop{\mathcal{H}\!\mathit{om}}\nolimits}
\def\SheafExt{\mathop{\mathcal{E}\!\mathit{xt}}\nolimits}
\def\Sch{\mathit{Sch}}
\def\Mor{\mathop{\mathrm{Mor}}\nolimits}
\def\Ob{\mathop{\mathrm{Ob}}\nolimits}
\def\Sh{\mathop{\mathit{Sh}}\nolimits}
\def\NL{\mathop{N\!L}\nolimits}
\def\CH{\mathop{\mathrm{CH}}\nolimits}
\def\proetale{{pro\text{-}\acute{e}tale}}
\def\etale{{\acute{e}tale}}
\def\QCoh{\mathit{QCoh}}
\def\Ker{\mathop{\mathrm{Ker}}}
\def\Im{\mathop{\mathrm{Im}}}
\def\Coker{\mathop{\mathrm{Coker}}}
\def\Coim{\mathop{\mathrm{Coim}}}

% Boxtimes
%
\DeclareMathSymbol{\boxtimes}{\mathbin}{AMSa}{"02}

%
% Macros for moduli stacks/spaces
%
\def\QCohstack{\mathcal{QC}\!\mathit{oh}}
\def\Cohstack{\mathcal{C}\!\mathit{oh}}
\def\Spacesstack{\mathcal{S}\!\mathit{paces}}
\def\Quotfunctor{\mathrm{Quot}}
\def\Hilbfunctor{\mathrm{Hilb}}
\def\Curvesstack{\mathcal{C}\!\mathit{urves}}
\def\Polarizedstack{\mathcal{P}\!\mathit{olarized}}
\def\Complexesstack{\mathcal{C}\!\mathit{omplexes}}
% \Pic is the operator that assigns to X its picard group, usage \Pic(X)
% \Picardstack_{X/B} denotes the Picard stack of X over B
% \Picardfunctor_{X/B} denotes the Picard functor of X over B
\def\Pic{\mathop{\mathrm{Pic}}\nolimits}
\def\Picardstack{\mathcal{P}\!\mathit{ic}}
\def\Picardfunctor{\mathrm{Pic}}
\def\Deformationcategory{\mathcal{D}\!\mathit{ef}}
\title{Derived completeness of the derived limit of a system killed by powers of the elements of an ideal}
\author{The Stacks Project Authors}
\date{stacks-project:more-algebra.tex:local-snapshot-20261012; GNU FDL 1.2 or later}
\begin{document}
\maketitle
\setcounter{subsection}{1}
\noindent\textit{Editorial note.} All delivered mathematical text is the author's own, copied from the pinned Stacks Project chapter file more-algebra.tex; no statement or proof step was rewritten, reordered or removed. The counted statement is the labeled lemma environment at source lines 26287-26294 (canonical tag 091W) and the counted proof is the proof environment at source lines 26296-26302, which immediately follows it and is its own proof.\par\medskip
\noindent\textit{Editorial note.} The proof uses the macro R-lim, the bold L of the derived tensor product and the torsion functor T of the shared preamble, all of which the retained preamble supplies verbatim.\par\medskip
\noindent\textit{Editorial note.} The two facts the proof invokes by name, namely that derived complete objects are closed under R-lim and that the vanishing of the torsion functor characterises derived completeness, are part of the framework established earlier in the same section; the counted proof states them as the reason for each reduction and contains no cross reference to a labelled item, so no navigation stub is needed.\par\medskip
\noindent\textit{Editorial note.} The chapter file pinned-original-source.tex carries only an input of the shared preamble on its first line, so no preamble text lives inside the chapter. The retained file sources/stacks-more-algebra/retained-preamble.tex is a mechanical, line-for-line copy of the project's single shared preamble as staged in the snapshot (Xy-pic setup plus the native theorem declarations and mathematical macros). Nothing was added, removed or redefined, and the excerpt serves only as preamble, never as a span coordinate.\par\medskip
\noindent\textit{Editorial note.} The source URL is the source-level Stacks Project URL, as required when more than one card is produced from one source; the per-claim tag is carried by claim\_key and by this note. The counted statement and proof cite no bibliography entry, so this card needs no bibliography.\par\medskip
\noindent\textit{Allowed standard prerequisites (editorial).}\begin{itemize}
\item The category of derived complete objects of D(A) with respect to an ideal I is closed under the derived limit of an inverse system. invoked at proof lines 26297-26298 as a standing property of derived complete objects. Reduces the assertion for the derived limit K' to the assertion for each of the objects K tensor\textasciicircum{}L K\_n.
\item The vanishing of the derived torsion functor T(-, f) on an object is equivalent to f being locally nilpotent on that object, so an object killed by a power of f is derived complete with respect to the principal ideal generated by f. used at proof lines 26299-26301, where the assumption that f\textasciicircum{}e is zero on K tensor\textasciicircum{}L K\_n is turned into T(K tensor\textasciicircum{}L K\_n, f) = 0. Shows that each term of the system is derived complete with respect to I, one generator at a time, which finishes the proof after the reduction.
\end{itemize}\medskip
\begin{lemma}
\label{lemma-naive-derived-completion}
Let $A$ be a ring and let $I \subset A$ be an ideal. Let $(K_n)$ be an inverse
system of objects of $D(A)$ such that for all $f \in I$ and $n$
there exists an $e = e(n, f)$ such that $f^e$ is zero on $K_n$.
Then for $K \in D(A)$ the object $K' = R\lim (K \otimes_A^\mathbf{L} K_n)$
is derived complete with respect to $I$.
\end{lemma}

\begin{proof}
Since the category of derived complete objects is preserved under $R\lim$
it suffices to show that each $K \otimes_A^\mathbf{L} K_n$ is derived
complete. By assumption for all $f \in I$ there is an $e$ such
that $f^e$ is zero on $K \otimes_A^\mathbf{L} K_n$. Of course this
implies that $T(K \otimes_A^\mathbf{L} K_n, f) = 0$ and we win.
\end{proof}

\end{document}
